% Introduction to the bishop vs bishop (opposite colours) chapter, for the 1500-1599 book only.
% The theory it needs comes from the library (lib/).
\section*{The defensive resources}
The defender should look for one of these. The attacker must make sure his opponent cannot use them.

\begin{itemize}
\item \textbf{First defensive system\index{Opposite-coloured bishops!First defensive system}}: the king blocks the passed pawns, the bishop protects the own pawn.

\item \textbf{Second defensive system\index{Opposite-coloured bishops!Second defensive system}}: the bishop watches the pawns along one diagonal, the king keeps the enemy king out.

\item \textbf{Bishop follows a pawn\index{Opposite-coloured bishops!Bishop follows a pawn}}: the bishop stays behind one pawn and watches it.

\item \textbf{Blocking on the colour of the opponent's bishop\index{Blocking on the colour of the opponent's bishop}}: pawns blocked there cannot move on.

\item \textbf{Switching the defenders\index{Switching the defenders}}: king and bishop swap their tasks.

\item \textbf{Removing the pawns\index{Exchanging or removing the pawns}}: give the bishop for the last pawns.
\end{itemize}

\section*{The attacking resources}
The attacker should look for one of these. The defender must make sure his opponent cannot use them.

\begin{itemize}
\item \textbf{King invasion\index{Opposite-coloured bishops!King invasion}}: the king walks in on the squares the enemy bishop cannot guard.

\item \textbf{Second passed pawn\index{Opposite-coloured bishops!Second passed pawn}}: two passed pawns far apart overload the defence.

\item \textbf{Two diagonals\index{Bishop!One-diagonal principle}}: a bishop that must watch two diagonals is torn apart.

\item \textbf{Breakthrough\index{Opposite-coloured bishops!Pawn breakthrough}}: a pawn sacrifice opens the way.

\item \textbf{Zugzwang\index{Zugzwang}}: the defender runs out of moves.
\end{itemize}

